\section{Conclusions}\label{sec:conclusions}

This study investigates the efficiency and limitations of optimised-multiplier error certification for discrete frictionless contact problems. The derived efficiency bound identifies the complementarity between exact contact forces and approximate gaps as a key factor governing certificate sharpness. For the stiffness-based inner product, the surrogate-minimising multiplier is characterised as the solution of a contact problem with a shifted obstacle. This interpretation yields a necessary and sufficient condition for exact multiplier recovery and a sufficient condition for exact error certification. Two distinct limitations are established: contact inconsistency in the displacement approximation can cause worst-case square-root degradation, whereas excluding a strictly active constraint from the multiplier support produces a positive asymptotic certificate floor. Numerical experiments support all tested theoretical results and demonstrate that the occurrence of exact certification depends on the correction strategy. These findings highlight the importance of preserving contact consistency on strictly active constraints and ensuring adequate multiplier support in certified reduced-order contact models.

The conditions for exact certification and the efficiency bounds involve quantities from the unknown full-order solution, limiting their direct use in online assessment. Future work should therefore focus on developing computable indicators of contact consistency and certificate sharpness. Alternative estimators incorporating contact-set information may also overcome the square-root degradation established for the considered certificate family. Further investigations are needed to characterise intermediate convergence rates, assess parameter-dependent stiffness matrices and larger contact problems, and evaluate certification costs relative to full-order solutions. Incorporating discretisation error into the certification framework remains another important direction.
